\documentclass[../main.tex]{subfiles}
\begin{document}

\begin{center}
    \Large{\textbf{ABSTRACT}}\\
\end{center}
\vspace{0.6cm}

\textbf{Background.} A model that reproduces how Vietnamese road users cross an intersection lets planners test an intersection design, a traffic-management measure, or an automated vehicle in simulation before doing so on the street. The standard behaviour models of microscopic simulators assume lane-disciplined, rule-compliant traffic, whereas motorcyclists in Vietnam run red lights one after another, start before the green, ride on the sidewalk in congestion, and turn left illegally by waiting inside the junction.

\textbf{Method.} The thesis separates rule generation from rule execution. A locally hosted large language model (LLM) generates and iteratively refines a behaviour rule specification made of vehicle-type parameters and conditional rules. Deterministic code validates the specification against a schema and a source-documented parameter space, executes it in the simulator while removing only the constraint that each violation needs removed, and measures its deviation from observations with ten metrics. The behaviour catalogue holds 11 behaviours drawn from the literature and from observation.

\textbf{Results.} The method is implemented in SUMO for the Le Van Luong -- Hoang Dao Thuy intersection in Hanoi during the morning peak, covering ten behaviours, with the local LLM \texttt{qwen3.5:9b}. The rule executor produces red-light running that default SUMO lacks (4.86\% versus 0\% of motorcycles), but the hand-written rule set still leaves vehicles waiting to enter the network and does not meet the acceptance criteria. The local LLM produces a valid specification within three attempts for 26 of 36 samples, and the validator catches all of its errors before any simulation run. In three campaigns with the same budget of 36 runs (Optuna and two LLM modes), no specification meets the acceptance criteria, so the LLM cannot yet be compared with Optuna.

\textbf{Limitations.} No video of the intersection exists yet. All results use substitute data built from published measurements at other intersections; they show that the workflow operates, not how closely the simulation matches real traffic.

\end{document}
